\section{Syntax}
\label{sec:syntax}

The syntax of \PurePy is given in \figref{syntax-stmt} and \figref{syntax-expr}. A \emph{block} is a non-empty sequence of
statements; either a top-level statement, a \pyinline{def} body, or a branch of a conditional. A
\emph{mutual region} $d_1 \ldots d_{n+1}$ is a maximal contiguous non-empty sequence of \pyinline{def}
statements in a block: any non-\pyinline{def} statement in that block, by definition, separates mutual
regions. Maximality is not enforced by the grammar in \figref{syntax-stmt}; the parser is responsible for
grouping consecutive \pyinline{def}s into a single region. The well-formedness rules treat each region as
one mutually-recursive binding (see \secref{mutual-recursion}).

\subsection{Lexical syntax}
\label{sec:lexical}

\todo{The grammar of \PurePy is an abstract syntax, and its metavariables range over values rather than
over source text: $n$ is a number and $w$ is a string, not the lexemes that denote them. How source text
becomes abstract syntax follows Python, including the indentation that delimits blocks; say where \PurePy
narrows those rules.}

\subsection{Contexts}
\label{sec:contexts}

Metavariables $x$, $f$, $c$ range over \emph{variables}; idiomatically $f$ is used when the variable is expected to denote a function and $c$ when it is expected to denote a class.

\begin{definition}[Context]
A context $\Gamma$ (or $\Delta$) is a finite map from variables to one of the following \emph{context entries} $\theta$:
\begin{itemize}
\item the variable's \emph{definite assignment} status $a \in \B = \set{\statusDefAssigned, \statusNotDefAssigned}$;
\item a \emph{module reference}, for a fully-qualified module name $q$: either a \emph{stub} $\modStub{q}$, or \emph{loaded} reference $\modLoaded{q}{\Gamma}$ carrying the \emph{signature} $\Gamma$ of $q$, the context specifying its members;
\item a \emph{class entry} $C$, of the form $\clsEntry{\Gamma}{q.c}{\seq{x}}{c'}$, recording the declaring context $\Gamma$, the qualified name $q.c$ of the class, sequence of distinct field names $\seq{x}$, and base class name $c'$ (or $\bot$ for no base).
\end{itemize}
\end{definition}

\noindent Module references arise from import statements. A module declares each class name at most once, so the qualified name of a class entry is the identity of the class. The definite-assignment terminology is borrowed from Java~\cite[\S16]{gosling2025}; $\statusDefAssigned$ indicates \emph{definitely assigned} and $\statusNotDefAssigned$ \emph{not definitely assigned} (or, equivalently, degenerate types ``any'' and ``none'').
It is useful to represent contexts as total functions on the set of all variables by extending the codomain with $\bot$ to indicate \emph{undefined}. We write $\Gamma, \bind{x}{\theta}$ for the context that sends $x$ to $\theta$ and otherwise agrees with $\Gamma$, define $\dom(\Gamma) = \set{x \mid \Gamma(x) \neq \bot}$, and write $\varnothing$ for the context where all variables are undefined. For a set of variables $X$, write $\statusDefAssigned(X)$ for the context $\set{\bind{x}{\statusDefAssigned} \mid x \in X}$; similarly $\statusNotDefAssigned(X)$.

\paragraph{Context operations.}

Contexts admit a non-commutative monoid of sequential composition $\Gamma \override \Delta$ (``override'') and
an idempotent commutative semigroup of parallel composition $\merge$ (``merge''). Together these form a
\emph{right near-semiring}~\cite{vanhoorn1967}: $\override$ right-distributes over $\merge$ (Lemma~1), but left
distributivity does not hold. Both operations are defined pointwise.
%
\begin{definition}[Context override]
{\small
\begin{align*}
\bot \override \theta &= \theta \override \bot = \theta \\
\theta \override \theta' &= \theta' \quad (\text{otherwise})
\end{align*}}
\end{definition}

\noindent Override is thus right-biased and has $\bot$ as both left and right unit.

\begin{definition}[Context merge]
For definite-assignment statuses $a, a' \in \B$:
{\small
\begin{align*}
\bot \merge \bot &= \bot \\
a \merge \bot = \bot \merge a &= \statusNotDefAssigned \\
a \merge a' &= a \meet a'
\end{align*}}
\end{definition}

\noindent Merge is otherwise undefined. Restricted to $\B$, it returns $\statusNotDefAssigned$ when exactly one argument is defined and the conjunction otherwise; so $\bot$ is not a unit.

\begin{definition}[Context extend]
\label{def:extend}
{\small
\begin{align*}
\modLoaded{q}{\Gamma} \extend \modLoaded{q}{\Gamma'} &= \modLoaded{q}{\Gamma \extend \Gamma'} \\
\modStub{q} \extend \modLoaded{q}{\Gamma} = \modLoaded{q}{\Gamma} \extend \modStub{q} &= \modLoaded{q}{\Gamma} \\
\bot \extend \theta = \theta \extend \bot &= \theta \\
\theta \extend \theta' &= \theta' \quad (\text{otherwise})
\end{align*}}
\end{definition}

\noindent Extend combines module references for the same module recursively, preferring loaded references to stubs, and is otherwise right-biased like override. It is used when processing import prefixes, where two imports rooted at the same package must combine rather than shadow.

\begin{lemma}[Right distributivity]
\[
(\Delta \merge \Delta') \override \Gamma = (\Delta \override \Gamma) \merge (\Delta' \override \Gamma).
\]
\end{lemma}

\noindent Left distributivity ($\Gamma \override (\Delta \merge \Delta') = (\Gamma \override \Delta) \merge
(\Gamma \override \Delta')$) does not hold. If it did, the following two programs would be equivalent, since
sequential composition with a preceding assignment could be distributed into the branches of a conditional.
But they are not: the left program is ill-formed, because a variable assigned before a conditional does not
remain definitely assigned if only some branches assign it (see also \secref{local-assignment}).
\vspace{-\medskipamount}

\begin{minipage}[t]{0.48\textwidth}
\begin{python}
def f(b):
    x = 1
    if b:
        x = 2
    else:
        pass
    return x  # error
\end{python}
\end{minipage}
\hfill
\begin{minipage}[t]{0.48\textwidth}
\begin{python}
def f(b):
    if b:
        x = 1
        x = 2
    else:
        x = 1
        pass
    return x  # ok
\end{python}
\end{minipage}

\begin{figure}
\small
\begin{center}
\begin{tabular}{rcll}
$s$ & $::=$ & $\exPass$ & (Pass) \\
    & $\mid$ & $\exAssign{x}{e}$ & (Assign) \\
    & $\mid$ & $e$ & (ExprStmt) \\
    & $\mid$ & $\exAssert{e}$ & (Assert) \\
    & $\mid$ & $\exAssertMsg{e}{e'}$ & (AssertMsg) \\
    & $\mid$ & $\exReturnNone$ & (ReturnNone) \\
    & $\mid$ & $\exReturn{e}$ & (Return) \\
    & $\mid$ & $\exIf{e}{b}{\exElif{e_1}{b_1}\;\ldots\;\exElif{e_n}{b_n}}$ & (If) \\
    & $\mid$ & $\exIfElse{e}{b}{\exElif{e_1}{b_1}\;\ldots\;\exElif{e_n}{b_n}}{b'}$ & (IfElse) \\
    & $\mid$ & $\exMatch{e}{\exCase{p_1}{b_1}\;\ldots\;\exCase{p_{n+1}}{b_{n+1}}}$ & (Match) \\
    & $\mid$ & $d_1 \;\ldots\; d_{n+1}$ & (MutualRegion) \\
    & $\mid$ & $\exDataclass{c}{\phi}$ & (Dataclass) \\
    & $\mid$ & $\exDataclassExtend{c}{c'}{\phi}$ & (DataclassExtend)
   \\[3mm]
$\iota$ & $::=$ & $\exImport{q}$ & (Import) \\
    & $\mid$ & $\exFromImport{q}{x_1, \ldots, x_{n+1}}$ & (FromImport)
   \\[3mm]
$\phi$ & $::=$ & $\exField{x_1} \;\ldots\; \exField{x_{n+1}}$ & (FieldList) \\
    & $\mid$ & $\pyinline{pass}$ & (EmptyFieldList)
   \\[3mm]
$b$ & $::=$ & $s_1 \;\ldots\; s_{n+1}$ & (Block)
   \\[3mm]
$t$ & $::=$ & $\iota_1 \;\ldots\; \iota_m \;\; s_1 \;\ldots\; s_n$ & (ModuleBody)
   \\[3mm]
$q$ & $::=$ & $x_1.x_2.\ldots.x_{n+1}$ & (Name)
   \\[3mm]
$d$ & $::=$ & $\exDef{f}{x_1, \ldots, x_n}{b}$ & (Def)
\end{tabular}
\end{center}
\caption{Modules and blocks}
\label{fig:syntax-stmt}
\end{figure}

\begin{figure}
\small
\centering
\ottgrammarfigure{
    \otte \ottinterrule
    \ottg \ottinterrule
    \ottl \ottinterrule
    \ottp \ottafterlastrule
}
\caption{Expressions and patterns}
\label{fig:syntax-expr}
\end{figure}

% \begin{center}
% \begin{tabular}{rcll}
% $e$ & $::=$ & $x$ & (Var) \\
%     & $\mid$ & $\ell$ & (Literal) \\
%     & $\mid$ & $\exApp{e}{e_1, \ldots, e_n}$ & (Call) \\
%     & $\mid$ & $\exConstruct{C}{e_1, \ldots, e_n}$ & (Construct) \\
%     & $\mid$ & $\exDict{x_1 \mathord{:} e_1, \ldots, x_n \mathord{:} e_n}$ & (Dict) \\
%     & $\mid$ & $\exProject{e}{x}$ & (Attribute) \\
%     & $\mid$ & $\exDProject{e}{e'}$ & (Subscript) \\
%     & $\mid$ & $\exBinaryApp{e}{\odot}{e'}$ & (BinOp) \\
%     & $\mid$ & $\exUnaryApp{\ominus}{e}$ & (UnaryOp) \\
%     & $\mid$ & $\exCondExpr{e_1}{e}{e_2}$ & (CondExpr) \\
%     & $\mid$ & $\exList{e_1, \ldots, e_n}$ & (List) \\
%     & $\mid$ & $\exTuple{e_1, \ldots, e_n}$ & (Tuple) \\
%     & $\mid$ & $\exListComp{e}{q_1\;\ldots\;q_n}$ & (ListComp) \\
%     & $\mid$ & $\exLambda{x_1, \ldots, x_n}{e}$ & (Lambda)
%    \\[3mm]
% $q$ & $::=$ & $\qualGuard{e}$ & (Guard) \\
%     & $\mid$ & $\qualGenerator{x}{e}$ & (Generator)
%    \\[3mm]
% $\ell$ & $::=$ & $n$ & (Number) \\
%     & $\mid$ & $\pyinline{True}$ & (True) \\
%     & $\mid$ & $\pyinline{False}$ & (False) \\
%     & $\mid$ & $\pyinline{None}$ & (None)
%    \\[3mm]
% $p$ & $::=$ & $\ell$ & (PatLiteral) \\
%     & $\mid$ & $-n$ & (PatNegLit) \\
%     & $\mid$ & $x$ & (PatVar) \\
%     & $\mid$ & $\patWild$ & (PatWildcard) \\
%     & $\mid$ & $\patList{p_1, \ldots, p_n}$ & (PatList) \\
%     & $\mid$ & $\patTuple{p_1, \ldots, p_n}$ & (PatTuple) \\
%     & $\mid$ & $\patClass{C}{p_1, \ldots, p_n, x_{n+1} \mathord{=} p_{n+1}, \ldots, x_{n+m} \mathord{=} p_{n+m}}$ & (PatClass) \\
%     & $\mid$ & $p\;\pyinline{as}\;x$ & (PatAs)
% \end{tabular}
% \end{center}


\paragraph{Operators.}

The metavariable $\odot$ ranges over the supported binary operators other than \pyinline{and} and
\pyinline{or}, and $\ominus$ over the supported unary
operators, both enumerated in \figref{precedence}. The operators \pyinline{and} and \pyinline{or} have their own
expression forms; they do not evaluate both operands. The grammar of \figref{syntax-expr} is intentionally ambiguous
on operator forms; disambiguation is by the precedence and associativity given in \figref{precedence},
which follow Python.

\begin{figure}
\small
\begin{center}
\begin{tabular}{rlll}
\hline
Level & Operators & Arity & Associativity \\
\hline
1 (lowest)  & \pyinline{or}                                                              & binary & left \\
2           & \pyinline{and}                                                             & binary & left \\
3           & \pyinline{not}                                                             & unary  & --- \\
4           & \pyinline{==}\quad\pyinline{!=}\quad\pyinline{<}\quad\pyinline{<=}\quad\pyinline{>}\quad\pyinline{>=} & binary & left \\
5           & \pyinline{+}\quad\pyinline{-}                                                & binary & left \\
6           & \pyinline{*}\quad\pyinline{/}\quad\pyinline{//}\quad\pyinline{\%}                & binary & left \\
7           & \pyinline{+}\quad\pyinline{-}                                                & unary  & --- \\
8 (highest) & \pyinline{**}                                                              & binary & right \\
\hline
\end{tabular}
\end{center}
\caption{Operator precedence (lowest binds loosest); follows Python.}
\label{fig:precedence}
\end{figure}


\subsection{Modules and programs}
\label{sec:modules-programs}

A \PurePy \emph{program} is a finite map $\mathcal{M}$ from \emph{fully-qualified module names}
(\figref{syntax-stmt}) to \emph{module bodies} $t$, of which one is a distinguished \emph{main module} named
$\pymath{\_\_main\_\_}$; here $\nameOf(q)$ is the string literal naming $q$, so
$\nameOf(\pymath{\_\_main\_\_}) = \pymath{"\_\_main\_\_"}$. We identify a module with its fully-qualified name, writing ``module $q$'' for the module named $q$. A module body is a sequence of imports
$\seq{\iota}$ followed by a sequence of statements $\seq{s}$; since imports are not statements, they can
appear neither in blocks (and hence not in function bodies) nor after the first non-import statement. The
well-formedness and evaluation rules check and evaluate the statements of each body $\mathcal{M}(q)$ with the
assignment $\exAssign{\pymath{\_\_name\_\_}}{\nameOf(q)}$ prepended, modelling the
$\pymath{\_\_name\_\_}$ binding the runtime supplies automatically. A second map $\mathcal{P}$ gives the
predefined modules of \figref{predefined}, which have no body: $\mathcal{P}$ takes each such name to its
environment (Definition~\ref{def:predefined}). The two domains are disjoint, and their union is
prefix-closed: the parent of any module in $\dom(\mathcal{M}) \cup \dom(\mathcal{P})$ is itself in
$\dom(\mathcal{M}) \cup \dom(\mathcal{P})$.\footnote{A source tree encodes such an $\mathcal{M}$: each module's body is its \texttt{.py}
file, each package's body is its \texttt{\_\_init\_\_.py}, and a package directory without an
\texttt{\_\_init\_\_.py} encodes the empty body (a namespace package, PEP 420, \url{https://peps.python.org/pep-0420/}). Prefix-closure thus holds by
construction.} The maps $\mathcal{M}$ and $\mathcal{P}$ are fixed and ambient in all
definitions that follow, as is the name $q$ of the enclosing module in the rules that check or evaluate a single module's body.


\begin{definition}[Predefined contexts and environments]
\label{def:predefined}
$\mathcal{P}$ takes each module $q$ of \figref{predefined} to an environment binding the members listed for
$q$, together with \pyinline{__name__}, to implementation-defined values. Its context is whatever types that
environment (\figref{well-formed-env}).
\end{definition}

\begin{figure}
\small
\centering
\begin{tabular}{lp{0.6\textwidth}}
\hline
\textbf{Module} & \textbf{Members} \\
\hline
\textup{\pyinline{builtins}} & \pyinline{print}, \pyinline{len}, \pyinline{range} \\
\pyinline{math} & \pyinline{pi}, \pyinline{e}, \pyinline{sqrt}, \pyinline{exp}, \pyinline{log}, \pyinline{sin}, \pyinline{cos}, \pyinline{tan}, \pyinline{floor}, \pyinline{ceil} \\
\pyinline{sys} & \pyinline{argv}, \pyinline{exit} \\
\pyinline{typing} & \pyinline{Any} \\
\pyinline{dataclasses} & \pyinline{dataclass} \\
\hline
\end{tabular}
\caption{Predefined modules and their members}
\label{fig:predefined}
\end{figure}

The metafunction $\assignsS$ extends by union from individual statements to blocks; on a module
body $\seq{\iota}\;\seq{s}$ it acts on $\seq{s}$. Write $q \prefixOf q'$ (\emph{prefix}) when $q' = q$ or $q' =
q.q''$ for some $q''$, and $\properPrefixOf$ for the strict order; the proper prefixes of a module name
are its \emph{ancestors}. The $\submods$ metafunction is
defined, with the class-related metafunctions, in \appendixref{sec:aux-modules}.

In this section, loading a module means checking its body and computing its signature
(\figref{well-formed-module}), the static counterpart of run-time loading (\secref{opsem}). Both
import forms load their target, together with its ancestors
(\figref{well-formed-imports}). A plain import also records each of these modules as a member of its
parent (the loads-as judgement), so that the binding created for the root of the target gives attribute
access along the whole chain. A from-import also loads any imported name that denotes a submodule. An
ancestor of the target is exempt from loading first if it is the importing module or one of its
ancestors, since at run time it may still be loading. The exemption applies only to from-imports, because
for a plain import
the only ancestor that could never load first is the importing module itself, and an import naming a proper
descendant of the importing module is excluded outright (\secref{no-load-dependent}). Loading a module
may thus require other modules to load first; the rules are interpreted inductively, so a program with an
import cycle is ill-formed, because a module in a cycle could finish loading only after itself. Program
well-formedness (\figref{well-formed-module}) requires the main module to be well-formed and hence,
recursively, every module the program can load; a module that is never imported is not checked.

\subsection{Well-formedness}

\subsubsection{Auxiliary definitions}

Three syntactic analyses support the well-formedness rules: $\assignsS(s)$, the set of variables assigned
anywhere in $s$ (without descending into nested function definitions, which have their own scope, so that
$\assignsS(s) = \dom(\Delta)$ when $s : \Assigns{\Delta}$); $\binds(p)$, the set of variables introduced by a
pattern $p$; and $\capturesS(s)$, the set of variables from the enclosing scope that are captured by closures
(lambdas or nested function definitions) within $s$, with companion $\capturesE(e)$ on expressions. All are
defined, by their principal cases, in \appendixref{sec:aux-syntactic}. $\capturesS$ extends to blocks and
module bodies by union (like $\assignsS$), with the singleton block identified with its statement.

\begin{definition}[Pattern subsumption]
The preorder $p \subsumed p'$ (\appendixref{sec:aux-subsumption}) holds when every value matched by $p$ is also matched by $p'$.
\end{definition}

\noindent Subsumption is used by the well-formedness rule for $\pyinline{match}$ to rule out unreachable cases and to check exhaustiveness; with the current pattern set, the latter reduces to requiring a catch-all (var or wildcard) case.

\subsubsection{Well-formedness rules}

A pattern list is well-formed (\figref{well-formed-pat}) when each pattern is linear (no repeated variable names) and no earlier pattern subsumes a later one. The \pyinline{match} well-formedness rules use this judgement to enforce reachability and linearity of case patterns.

\begin{figure}
\flushleft \shadebox{$\Gamma \vdash p_1 \cons \ldots \cons p_n$}
\begin{mathpar}
\small
\inferrule*[lab={\ruleName{pat-lit}}]
{ \strut }
{ \Gamma \vdash \ell }
\and
\inferrule*[lab={\ruleName{pat-neg-lit}}]
{ \strut }
{ \Gamma \vdash -n }
\and
\inferrule*[lab={\ruleName{pat-var}}]
{ \strut }
{ \Gamma \vdash x }
\and
\inferrule*[lab={\ruleName{pat-wild}}]
{ \strut }
{ \Gamma \vdash \patWild }
\and
\inferrule*[lab={\ruleName{pat-list}}]
{ \Gamma \vdash p_i \quad (\forall i) }
{ \Gamma \vdash \patList{p_1, \ldots, p_n} }
\and
\inferrule*[lab={\ruleName{pat-tuple}}]
{ \Gamma \vdash p_i \quad (\forall i) }
{ \Gamma \vdash \patTuple{p_1, \ldots, p_n} }
\and
\inferrule*[lab={\ruleName{pat-as}}]
{ \Gamma \vdash p }
{ \Gamma \vdash p\;\pyinline{as}\;x }
\and
\inferrule*[lab={\ruleName{pat-class}}]
{
  \Gamma \vdash q \names d(x_1, \ldots, x_n, x_{n+1}, \ldots, x_{n+m}) \\
  \set{y_1, \ldots, y_m} = \set{x_{n+1}, \ldots, x_{n+m}} \\
  \Gamma \vdash p_i \quad (\forall i)
}
{
  \Gamma \vdash \patClass{q}{p_1, \ldots, p_n, y_1 \mathord{=} p_{n+1}, \ldots, y_m \mathord{=} p_{n+m}}
}
\and
\inferrule*[lab={\ruleName{cases-nil}}]
{ \strut }
{ \Gamma \vdash \varepsilon }
\and
\inferrule*[lab={\ruleName{cases-cons}}]
{
  \Gamma \vdash p_1 \\
  \Gamma \vdash p_2 \cons \ldots \cons p_n \\
  \text{variables in } p_1 \text{ distinct} \\
  p_i \not\subsumed p_1 \quad (\forall i.\, 2 \le i \le n)
}
{
  \Gamma \vdash p_1 \cons p_2 \cons \ldots \cons p_n
}
\end{mathpar}
\caption{Case patterns are well-formed under $\Gamma$}
\label{fig:well-formed-pat}
\end{figure}


Every well-formed statement either \emph{definitely returns}, or \emph{definitely assigns} a context
$\Delta$, which may be empty. The base cases are that \pyinline{return e} definitely returns, and \pyinline{x
= e} definitely assigns $\set{\bind{\pyinline{x}}{\statusDefAssigned}}$. \emph{Result types} $R$ classify these outcomes: $R$ is either $\Returns$ (definitely returns) or $\Assigns{\Delta}$ (definitely assigns $\Delta$). Block well-formedness is given in
\figref{well-formed}; expression well-formedness in \figref{well-formed-expr}.

Context operations $\override$ and $\merge$ lift to a right semiring over result types $R$, with $\Assigns{\varnothing}$ as left and right unit for $\override$ and $\Returns$ as unit for $\merge$ and zero for $\override$.

\begin{definition}[Override and merge on result types]
{\small
\begin{align*}
\Assigns{\Delta} \merge \Assigns{\Delta'} & = \Assigns{(\Delta \merge \Delta')} \\
R \merge \Returns = \Returns \merge R & = R \\
\\
\Assigns{\Delta} \override \Assigns{\Delta'} & = \Assigns{(\Delta \override \Delta')} \\
R \override \Returns = \Returns \override R & = \Returns
\end{align*}}
\end{definition}

\noindent The $\Returns \override R$ case never arises because the well-formedness judgement excludes unreachable statements (\secref{unreachable}).

\begin{lemma}[Right distributivity (result types)]
\[
(R \merge R') \override S = (R \override S) \merge (R' \override S)
\]
\end{lemma}

\noindent A module attribute reference is well-formed only when it resolves to a definitely assigned member
(rule \ruleName{attr-module}); see \secref{no-load-dependent}.

\begin{figure}
\flushleft \shadebox{$\Gamma; \Lambda \vdash b : R$}
\begin{mathpar}
\small
\inferrule*[lab={\ruleName{pass}}]
{
  \strut
}
{
  \Gamma; \Lambda \vdash \exPass : \Assigns{\varnothing}
}
\and
\inferrule*[lab={\ruleName{assign}}]
{
  \Gamma; \Lambda \vdash e \\
  x \notin \capturesE(e)
}
{
  \Gamma; \Lambda \vdash \exAssign{x}{e} : \Assigns{\set{x : \statusDefAssigned}}
}
\and
\inferrule*[lab={\ruleName{expr-stmt}}]
{
  \Gamma; \Lambda \vdash e
}
{
  \Gamma; \Lambda \vdash e : \Assigns{\varnothing}
}
\and
\inferrule*[lab={\ruleName{assert}}]
{
  \Gamma; \Lambda \vdash e
}
{
  \Gamma; \Lambda \vdash \exAssert{e} : \Assigns{\varnothing}
}
\and
\inferrule*[lab={\ruleName{assert-msg}}]
{
  \Gamma; \Lambda \vdash e \\
  \Gamma; \Lambda \vdash e'
}
{
  \Gamma; \Lambda \vdash \exAssertMsg{e}{e'} : \Assigns{\varnothing}
}
\and
\inferrule*[lab={\ruleName{return-none}}]
{
  \strut
}
{
  \Gamma; \Lambda \vdash \exReturnNone : \Returns
}
\and
\inferrule*[lab={\ruleName{return}}]
{
  \Gamma; \Lambda \vdash e
}
{
  \Gamma; \Lambda \vdash \exReturn{e} : \Returns
}
\and
\inferrule*[lab={\ruleName{if}}]
{
  \Gamma; \Lambda \vdash e_i \quad (\forall i) \\
  \Gamma; \Lambda \vdash b_i : R_i \quad (\forall i)
}
{
  \Gamma; \Lambda \vdash \exIf{e_1}{b_1}{\exElif{e_2}{b_2}\;\ldots\;\exElif{e_{n+1}}{b_{n+1}}} : R_1 \merge \ldots R_{n + 1} \merge \ldots \merge \Assigns{\varnothing}
}
\and
\inferrule*[lab={\ruleName{if-else}}]
{
  \Gamma; \Lambda \vdash e_i \quad (\forall i) \\
  \Gamma; \Lambda \vdash b_i : R_i \quad (\forall i)
}
{
  \Gamma; \Lambda \vdash \exIfElse{e_1}{b_1}{\exElif{e_2}{b_2}\;\ldots\;\exElif{e_{n+1}}{b_{n+1}}}{b_{n+2}} : R_1 \merge \ldots \merge R_{n+2}
}
\and
\inferrule*[lab={\ruleName{mutual}}]
{
  \Gamma \runion \statusDefAssigned(\seq{f}) \runion \statusDefAssigned(\seq{x}_{i}) \runion
  \statusNotDefAssigned(\assignsS(b_i) \setminus \seq{x}_i); \Lambda \vdash b_i : R_i
  \quad (\forall i)
  \\
  f_1, \ldots, f_{n+1} \text{ distinct}
}
{
  \Gamma; \Lambda \vdash \exDef{f_1}{\vec{x}_1}{b_1} \cons \ldots \cons \exDef{f_{n+1}}{\vec{x}_{n+1}}{b_{n+1}} : \Assigns{\statusDefAssigned(\seq{f})}
}
\and
\inferrule*[lab={\ruleName{match}}]
{
  \Gamma; \Lambda \vdash e \\
  \Gamma \runion \statusDefAssigned(\binds(p_i)); \Lambda \vdash b_i : R_i \quad (\forall i) \\
  S_i = \Assigns{\statusDefAssigned(\binds(p_i))} \runion R_i \quad (\forall i) \\
  \Lambda \vdash p_1 \cons \ldots \cons p_{n+1} \\
  p_{n+1} \in \set{x, \patWild}
}
{
  \Gamma; \Lambda \vdash \exMatch{e}{\exCase{p_1}{b_1}\;\ldots\;\exCase{p_{n+1}}{b_{n+1}}} : S_1 \merge \ldots \merge S_{n+1}
}
\and
\inferrule*[lab={\ruleName{match-partial}}]
{
  \Gamma; \Lambda \vdash e \\
  \Gamma \runion \statusDefAssigned(\binds(p_i)); \Lambda \vdash b_i : R_i \quad (\forall i) \\
  S_i = \Assigns{\statusDefAssigned(\binds(p_i))} \runion R_i \quad (\forall i) \\
  \Lambda \vdash p_1 \cons \ldots \cons p_{n+1} \\
  p_{n+1} \notin \set{x, \patWild}
}
{
  \Gamma; \Lambda \vdash \exMatch{e}{\exCase{p_1}{b_1}\;\ldots\;\exCase{p_{n+1}}{b_{n+1}}} : S_1 \merge \ldots \merge S_{n+1} \merge \Assigns{\varnothing}
}
\and
\inferrule*[lab={\ruleName{cons}}]
{
  \Gamma; \Lambda \vdash s : \Assigns{\Delta} \\
  \Gamma \runion \Delta; \Lambda \vdash b : R \\
  \capturesS(s) \cap \assignsS(b) = \varnothing
}
{
  \Gamma; \Lambda \vdash s \cons b : \Assigns{\Delta} \runion R
}
\end{mathpar}
\caption{Block $b$ is well-formed in $\Gamma$ and $\Lambda$}
\label{fig:well-formed}
\end{figure}

\begin{figure}
\begin{subfigure}{\textwidth}
\flushleft \shadebox{$\Gamma \vdash q \names \theta$}
\begin{mathpar}
\small
\inferrule*[lab={\ruleName{simple-module}}]
{
  x : q \in \Gamma
}
{
  \Gamma \vdash x \names q
}
\and
\inferrule*[lab={\ruleName{simple-class}}]
{
  c : (q', \vec{y}, c') \in \Gamma \\
  \fields(q'.c) = \vec{x}
}
{
  \Gamma \vdash c \names q'.c(\vec{x})
}
\and
\inferrule*[lab={\ruleName{qualified-module}}]
{
  \Gamma \vdash q \names q' \\
  q'.x \in \dom(\mathcal{M})
}
{
  \Gamma \vdash q.x \names q'.x
}
\and
\inferrule*[lab={\ruleName{qualified-class}}]
{
  \Gamma \vdash q \names q' \\
  \Gamma_{q'}(c) = (q', \vec{y}, c'') \\
  \fields(q'.c) = \vec{x}
}
{
  \Gamma \vdash q.c \names q'.c(\vec{x})
}
\end{mathpar}
\subcaption{Name $q$ resolves to context entry $\theta$ under $\Gamma$}
\end{subfigure}

\vspace{2mm}

\begin{subfigure}{\textwidth}
\flushleft \shadebox{$\Gamma \vdash e$}
\begin{mathpar}
\small
\inferrule*[lab={\ruleName{var}}]
{
  \Gamma(x) = \statusDefAssigned
}
{
  \Gamma \vdash x
}
\and
\inferrule*[lab={\ruleName{lambda}}]
{
  \Gamma \runion \statusDefAssigned(\seq{x}) \vdash e
}
{
  \Gamma \vdash \exLambda{\seq{x}}{e}
}
\and
\inferrule*[lab={\ruleName{constr}}]
{
  \Gamma \vdash q \names d(x_1, \ldots, x_n) \\
  \Gamma \vdash e_i \quad (\forall i)
}
{
  \Gamma \vdash \exConstruct{q}{e_1, \ldots, e_n}
}
\and
\inferrule*[lab={\ruleName{attr-module}}]
{
  \Gamma \vdash e \names q \\
  y \in \dom(\Gamma_q)
}
{
  \Gamma \vdash e.y
}
\and
\inferrule*[lab={\ruleName{attr-object}}]
{
  \Gamma \vdash e
}
{
  \Gamma \vdash e.y
}
\end{mathpar}
\subcaption{Expression $e$ is well-formed under $\Gamma$}
\end{subfigure}
\caption{Expression well-formedness}
\label{fig:well-formed-expr}
\end{figure}

\begin{figure}
\flushleft \shadebox{$\Gamma \vdash m : \Delta$}
\begin{mathpar}
\small
\inferrule*[lab={\ruleName{empty}}]
{
  \strut
}
{
  \Gamma \vdash \varepsilon : \varnothing
}
\and
\inferrule*[lab={\ruleName{import}}]
{
  M' = x_1.x_2.\ldots.x_{n+1} \\
  \Gamma \runion \statusDefAssigned(\set{x_1}) \vdash m : \Delta
}
{
  \Gamma \vdash \exImport{M'} \cons m : \statusDefAssigned(\set{x_1}) \runion \Delta
}
\and
\inferrule*[lab={\ruleName{from-import}}]
{
  \Gamma \runion \statusDefAssigned(\set{x_1, \ldots, x_{n+1}}) \vdash m : \Delta
}
{
  \Gamma \vdash \exFromImport{M'}{x_1, \ldots, x_{n+1}} \cons m : \statusDefAssigned(\set{x_1, \ldots, x_{n+1}}) \runion \Delta
}
\and
\inferrule*[lab={\ruleName{statement}}]
{
  \Gamma \vdash s : \Assigns{\Delta} \\
  \Gamma \runion \Delta \vdash m : \Delta' \\
  \capturesS(s) \cap \assignsS(m) = \varnothing
}
{
  \Gamma \vdash s \cons m : \Delta \runion \Delta'
}
\and
\inferrule*[lab={\ruleName{class}}]
{
  \decls(\phi) = x_1, \ldots, x_n \\
  x_1, \ldots, x_n \text{ distinct} \\
  \Gamma \vdash m : \Delta
}
{
  \Gamma \vdash \exDataclass{C}{\phi} \cons m : \Delta
}
\and
\inferrule*[lab={\ruleName{class-extend}}]
{
  B \in \dom(\Lambda_M) \\
  \decls(\phi) = x_1, \ldots, x_n \\
  x_1, \ldots, x_n \text{ distinct and disjoint from } \fields(M.B) \\
  \Gamma \vdash m : \Delta
}
{
  \Gamma \vdash \exDataclassExtend{C}{B}{\phi} \cons m : \Delta
}
\end{mathpar}
\caption{Body $m$ of module $M$ is well-formed in $\Gamma$}
\label{fig:well-formed-module}
\end{figure}


The signature $\Gamma$ a module's load judgement produces (\figref{well-formed-module}) comprises the classes
and variables a module defines, together with its intrinsic \pyinline{__name__} and a stub for each submodule,
never the names it imports: following PEP 484\footnote{\url{https://peps.python.org/pep-0484/}}, an imported
name is a private dependency, not re-exported as part of the interface. Submodule stubs and the module's own
bindings are disjoint, because the module rules exclude a module that binds a name that is also the name of
a submodule (\secref{no-load-dependent}).

\subsection{Examples}

\subsubsection{Assignment as definition}
\label{sec:assignment-as-definition}

Although \PurePy is pure and all variable usage is well-scoped, local variables are introduced using an
\emph{assignment statement}, for compatibility with Python. If a variable is definitely assigned in any branch
of a conditional, then for that variable to be usable after the conditional, every other branch of the
conditional must either definitely return, or definitely assign that variable too. Thus the following are both
well-formed:

\begin{figure}[!ht]
\centering
\begin{subfigure}{0.48\textwidth}
\centering
\begin{python}
if a < b:
   comp = "smaller"
else:
   comp = "bigger"
return comp
\end{python}
\end{subfigure}
\hfill
\begin{subfigure}{0.48\textwidth}
\centering
\begin{python}
if a < b:
   return "smaller"
else:
   comp = "bigger"
return comp
\end{python}
\end{subfigure}
\end{figure}

\noindent The definite assignment set for a conditional that has at least one branch which definitely assigns
is calculated by merging the definite assignment sets for all branches which definitely assign, using the
$\merge$ operator on environments.

Definite assignment thus gives each reference a unique lexical referent: a variable is introduced by
an assignment, possibly jointly by the branches of a conditional, and a name may be referenced only where one
such introduction covers every path. This is stricter than avoiding Python's \pyinline{NameError}: the
left-hand program of \secref{contexts} assigns \pyinline{x} on every run, but is rejected because the
reference would denote the branch-local variable on one path and the variable it hides on the other.

\subsubsection{Local assignment in conditional clauses}
\label{sec:local-assignment}

For flexibility, we permit \emph{branch-local} variables. The assignment of \pyinline{x} below (left) is
permitted, and can be used locally in that branch of the conditional:

\begin{figure}[!ht]
\centering
\begin{subfigure}{0.48\textwidth}
\centering
\begin{python}
def foo():
  z = 6
  if b:
    x = 3
    y = x + 1
  else:
    y = 0
  return z
\end{python}
\caption{Branch-local definition of \pyinline{x}}
\end{subfigure}
\hfill
\begin{subfigure}{0.48\textwidth}
\centering
\begin{python}
def foo():
  x = 6
  if b:
    x = 3
    y = x + 1
  else:
    y = 0
  return x # error: partially defined
\end{python}
\caption{Now with use of \pyinline{x} after the conditional}
\end{subfigure}
\end{figure}

\noindent However such variables only appear to be branch-local. The definition of $\merge$ actually ensures
that a variable assigned in any branch is considered assigned by the conditional; however unless it is
assigned in \emph{every} branch that does not definitely return, it will be bound to $\statusNotDefAssigned$ in the
corresponding environment. Any attempt to refer to such a variable is illegal, as it refers to a variable
which may not be assigned. (Thus the right-hand program is ill-formed, despite the enclosing \pyinline{x = 6},
because in the \pyinline{return x} statement, the inner assignment shadows the the outer one, and the inner
assignment is only partial.) If we were to allow the program on the right, there would be no straightforward
way to assign it a semantics which is both pure, and in agreement with Python, since in Python the function
returns \pyinline{3} in the case where \pyinline{b} is \pyinline{True}.

\subsubsection{Function-level scoping of local variables}

In Python, all local variables are scoped to the entire function body, not just from the point of
assignment onwards. The $\assignsS$ function in the \ruleName{def} rule reflects this: all variables
assigned anywhere in the function body are introduced at $\statusNotDefAssigned$ before the body is checked. This means
the following program is statically rejected by \PurePy, whereas Python rejects it at runtime with an
\pyinline{UnboundLocalError}:

\begin{python}
y = 7

def f():
    x = y  # error: y is local (due to assignment below)
    y = 8

f()
\end{python}

\noindent Because \pyinline{y} appears on the left-hand side of an assignment in \pyinline{f}, it is in
$\assignsS(b)$ for the body of \pyinline{f}, and is therefore introduced at $\statusNotDefAssigned$ at the top of the
function. The reference to \pyinline{y} in \pyinline{x = y} then fails the well-formedness check, since
$\Gamma(\pyinline{y}) = \statusNotDefAssigned$ rather than $\Gamma(\pyinline{y}) = \statusDefAssigned$.

\subsubsection{Reassignment of captured variables}
\label{sec:captured-vars}

The exclusion of captured-variable reassignment (\secref{no-mutable-state}) is enforced by a side condition
on the \ruleName{cons} rule: no variable captured by a statement $s$ may be reassigned in the rest of the block
$b$. In the example of \secref{no-mutable-state}, \pyinline{x} is captured by the definition of
\pyinline{g} and reassigned by \pyinline{x = 6}, so the program is rejected. A dataclass declaration counts as
an assignment of its class name for these definitions, though re-declaring a class name is excluded outright
(\secref{no-first-class}).

The same issue arises when a variable is captured and reassigned in the \emph{same} statement:

\begin{python}
x = 5
def x():
    return x
x()  # Python: returns the function (late binding); pure: returns 5
\end{python}

\noindent The side condition on the \ruleName{assign} rule rejects programs like this, requiring that the
variable being bound is not captured by the expression. The \ruleName{def} rule, by contrast, binds
each $f_i$ in the typing context for every body $b_i$, so the function names may legitimately appear in
$\capturesS$ of those bodies. This is what permits self-recursion (the singleton-region case) and
mutual recursion within a region.

\subsubsection{Mutual recursion}
\label{sec:mutual-recursion}

The \ruleName{def} rule binds the names $f_1, \ldots, f_{n+1}$ of a mutual region $d_1 \ldots d_{n+1}$
simultaneously. Each body $b_i$ is checked in a context that includes all the $f_j$ as definitely
assigned, so any $f_j$ may reference any other $f_k$ within the region (the singleton case $n = 0$
permits self-recursion). The names in a region must be distinct.

Inserting any non-\pyinline{def} statement between two \pyinline{def}s splits them into separate regions,
breaking the mutual reference:

\begin{python}
def even(n):
    if n == 0:
        return True
    return odd(n - 1)

x = even(0)  # intervening statement separates the regions

def odd(n):
    if n == 0:
        return False
    return even(n - 1)
\end{python}

\noindent This program is rejected statically: \pyinline{even} forms a singleton region in which
\pyinline{odd} is not yet bound, so the reference to \pyinline{odd} in its body fails the
\ruleName{var} rule. Python accepts it: the call \pyinline{even(0)} does not actually reach
\pyinline{odd}, and Python's late binding means \pyinline{odd} is only looked up when needed.

Each branch of a conditional may contain its own mutual regions, subject to the usual definite assignment
rules. The following is well-formed because both branches define both names:

\begin{python}
b = True

if b:
    def f():
        return g()
    def g():
        return "via mutual region"
else:
    def f():
        return "f independent"
    def g():
        return "g independent"

print(f())
print(g())
\end{python}

A region containing two definitions of the same name is rejected directly by the \ruleName{def}
rule's distinct-names side condition:

\begin{python}
def g():
    return 0

def f():
    return g()
def g():     # error: g defined twice in the same region
    return 1
\end{python}

\noindent (Python accepts this; late binding makes \pyinline{f}'s call to \pyinline{g} resolve to the
second \pyinline{def g}.)

Wrapping the second \pyinline{def f; def g} pair in a conditional sidesteps the duplicate-name issue:
the outer \pyinline{def g} and the inner \pyinline{def g} are now in different blocks, so they no longer
fall in the same region. One might still expect a captured-variable problem, since the inner
\pyinline{def g} could be seen as rebinding the outer \pyinline{g} that \pyinline{f} captures. No such
problem arises:

\begin{python}
def g():
    return 0

if b:
    def f():
        return g()
    def g():
        return 1
\end{python}

\noindent The inner \pyinline{def f; def g} form a branch-local mutual region, so \pyinline{f}'s body
resolves \pyinline{g} to its sibling rather than to the outer \pyinline{g}. The two are bound
simultaneously, so \pyinline{f} cannot be called before the inner \pyinline{g} exists, and the value
\pyinline{f} sees is the same as the value Python's late binding finds at call time.

\subsubsection{Pattern variables escape the match}

A variable bound by a \pyinline{case} pattern is in scope after the \pyinline{match} statement, just as
a variable assigned inside a conditional branch is in scope after the conditional. Python treats matching
this way; \PurePy follows.

\begin{python}
v = 5
match v:
    case x:
        pass
print(x)    # prints 5
\end{python}

Here the catch-all \pyinline{case x:} binds \pyinline{x} in the body's scope, and the binding survives into
the enclosing scope. Definite assignment after the match follows the same merge logic as a conditional: if
every case binds the variable (or returns), it is definitely assigned afterwards; otherwise it is not (see
also \secref{local-assignment}).

\subsubsection{Unreachable code}
\label{sec:unreachable}
A non-empty sequence of statements returns only if the last statement returns, and all preceding statements
assign. Thus the following is illegal:

\begin{python}
def foo(x):
  return x + 1
  return x    # error: unreachable
\end{python}

\begin{DeletedBlock}
\subsubsection{Conditionals without \pyinline{else} clauses}

An \pyinline{if} statement without an \pyinline{else} clause is equivalent to an \pyinline{if} statement with
an \pyinline{else} clause whose body is a $\exPass$ statement; such a statement never returns.

\subsubsection{Implicit \pyinline{None}}

A \pyinline{return} statement without an expression is equivalent to \pyinline{return None}.

\subsubsection{Implicit \pyinline{return}}

A function whose body $s$ assigns rather than returns is equivalent to a function whose body is $s$ followed
by \pyinline{return None}. This ensures every function returns, even if it is not explicitly specified.

\subsubsection{\pyinline{assert} without a message}

An \pyinline{assert e} statement without a message is equivalent to \pyinline{assert e, ""}.
\end{DeletedBlock}
